% GENERATED FILE. Edit canonical lessons, then run npm run generate:books.
% canonical-source-hash: fnv1a64:fcf13dc80bcdc769
% canonical-lessons: TE-C373-atmakatha, TE-C373-jivitacaritra, TE-C373-sameta, TE-C373-burrakatha, TE-C373-srikaram-cuttu, TE-R373-review-366-373

\chapter{Lives in Writing}
\label{ch:te-lives-in-writing}
\hlchaptermodality{373}

\begin{chapteropening}
\textbf{By the end of this chapter:} \emph{I can say an autobiography, a biography, a proverb, a burrakatha ballad and to make a start.}

Complete the last lesson of chapter 373: I can say an autobiography, a biography, a proverb, a burrakatha ballad and to make a start.
\end{chapteropening}

\section[\texorpdfstring{\te{ఆత్మకథ} (ātmakatha)}{ఆత్మకథ (ātmakatha)}]{\te{ఆత్మకథ} (ātmakatha) — an autobiography}

\label{lesson:TE-C373-atmakatha}



\begin{quote}
Before the new one: say the Telugu for imagination, then the Telugu for a made-up story.
\end{quote}

\subsection*{You'll want to know: \te{ఆత్మకథ}}
\textbf{\te{ఆత్మకథ}} — \emph{ātmakatha} — \textquotedblleft{}an autobiography\textquotedblright{}.

\begin{quote}
\te{ఆ} \te{నాయకుడు} \te{తన} \te{ఆత్మకథ} \te{రాశాడు}.
\end{quote}

\begin{quote}
\emph{ā nāyakuḍu tana ātmakatha rāśāḍu.}
\end{quote}

\begin{quote}
That leader wrote his autobiography.
\end{quote}

\te{ఆత్మ}, self, and \te{కథ}: one's own story.

\begin{cousinweb}
\textbf{\te{ఆత్మకథ}} (\emph{ātmakatha}) is built from Sanskrit \emph{ātman}, \textquotedblleft{}breath; self\textquotedblright{} (CDIAL 1135), and \emph{kathā}, \textquotedblleft{}a story\textquotedblright{}.
\end{cousinweb}

\begin{grammarlens}[title={What you've built}]
One more word for telling what happened.
\end{grammarlens}

\subsection*{Your turn}
\begin{itemize}
\raggedright
  \item \emph{Say it:} \emph{ātmakatha}
  \item \emph{Say it:} \emph{ātmakatha}, once more
  \item \emph{Say it:} \emph{ā nāyakuḍu tana ātmakatha rāśāḍu.} — That leader wrote his autobiography.
  \item \emph{Recall:} say the Telugu for imagination, then the Telugu for a made-up story, then say \emph{ātmakatha} again
\end{itemize}

\begin{tcolorbox}[breakable,colback=teal!4,colframe=teal!35!black,title={Before you move on}]
What does \te{ఆత్మకథ} mean? (\textquotedblleft{}An autobiography\textquotedblright{}.) Say it once more.
\end{tcolorbox}
\section[\texorpdfstring{\te{జీవితచరిత్ర} (jīvitacaritra)}{జీవితచరిత్ర (jīvitacaritra)}]{\te{జీవితచరిత్ర} (jīvitacaritra) — a biography}

\label{lesson:TE-C373-jivitacaritra}



\begin{quote}
Before the new one: say the Telugu for a made-up story, then the Telugu for an autobiography.
\end{quote}

\subsection*{You'll want to know: \te{జీవితచరిత్ర}}
\textbf{\te{జీవితచరిత్ర}} — \emph{jīvitacaritra} — \textquotedblleft{}a biography\textquotedblright{}.

\begin{quote}
\te{పిల్లలు} \te{ఆ} \te{వీరుడి} \te{జీవితచరిత్ర} \te{చదివారు}.
\end{quote}

\begin{quote}
\emph{pillalu ā vīruḍi jīvitacaritra cadivāru.}
\end{quote}

\begin{quote}
The children read the hero's biography.
\end{quote}

\te{జీవితం}, a life, and \te{చరిత్ర}, a history.

\begin{cousinweb}
\textbf{\te{జీవితచరిత్ర}} (\emph{jīvitacaritra}) is built from Sanskrit \emph{jīvita}, \textquotedblleft{}living, life\textquotedblright{} (CDIAL 5252), and \emph{caritra}, \textquotedblleft{}behaviour, conduct\textquotedblright{} (CDIAL 4691).
\end{cousinweb}

\begin{grammarlens}[title={What you've built}]
One more word for telling what happened.
\end{grammarlens}

\subsection*{Your turn}
\begin{itemize}
\raggedright
  \item \emph{Say it:} \emph{jīvitacaritra}
  \item \emph{Say it:} \emph{jīvitacaritra}, once more
  \item \emph{Say it:} \emph{pillalu ā vīruḍi jīvitacaritra cadivāru.} — The children read the hero's biography.
  \item \emph{Recall:} say the Telugu for a made-up story, then the Telugu for an autobiography, then say \emph{jīvitacaritra} again
\end{itemize}

\begin{tcolorbox}[breakable,colback=teal!4,colframe=teal!35!black,title={Before you move on}]
What does \te{జీవితచరిత్ర} mean? (\textquotedblleft{}A biography\textquotedblright{}.) Say it once more.
\end{tcolorbox}
\section[\texorpdfstring{\te{సామెత} (sāmeta)}{సామెత (sāmeta)}]{\te{సామెత} (sāmeta) — a proverb; a saying}

\label{lesson:TE-C373-sameta}



\begin{quote}
Before the new one: say the Telugu for an autobiography, then the Telugu for a biography.
\end{quote}

\subsection*{You'll want to know: \te{సామెత}}
\textbf{\te{సామెత}} — \emph{sāmeta} — \textquotedblleft{}a proverb; a saying\textquotedblright{}.

\begin{quote}
\te{పెద్దలు} \te{ఒక} \te{సామెత} \te{చెప్పారు}.
\end{quote}

\begin{quote}
\emph{peddalu oka sāmeta ceppāru.}
\end{quote}

\begin{quote}
The elders quoted a proverb.
\end{quote}

Plural: \te{సామెతలు}, \emph{sāmetalu}.

Telugu is rich in proverbs; elders drop them into talk.

\begin{cousinweb}
\textbf{\te{సామెత}} (\emph{sāmeta}) is the Telugu word for a proverb; its origin is not settled.
\end{cousinweb}

\begin{grammarlens}[title={What you've built}]
One more word for telling what happened.
\end{grammarlens}

\subsection*{Your turn}
\begin{itemize}
\raggedright
  \item \emph{Say it:} \emph{sāmeta}
  \item \emph{Say it:} \emph{sāmeta}, once more
  \item \emph{Say it:} \emph{peddalu oka sāmeta ceppāru.} — The elders quoted a proverb.
  \item \emph{Recall:} say the Telugu for an autobiography, then the Telugu for a biography, then say \emph{sāmeta} again
\end{itemize}

\begin{tcolorbox}[breakable,colback=teal!4,colframe=teal!35!black,title={Before you move on}]
What does \te{సామెత} mean? (\textquotedblleft{}A proverb; a saying\textquotedblright{}.) Say it once more.
\end{tcolorbox}
\section[\texorpdfstring{\te{బుర్రకథ} (burrakatha)}{బుర్రకథ (burrakatha)}]{\te{బుర్రకథ} (burrakatha) — a burrakatha ballad}

\label{lesson:TE-C373-burrakatha}



\begin{quote}
Before the new one: say the Telugu for a biography, then the Telugu for a proverb.
\end{quote}

\subsection*{You'll want to know: \te{బుర్రకథ}}
\textbf{\te{బుర్రకథ}} — \emph{burrakatha} — \textquotedblleft{}a burrakatha ballad\textquotedblright{}.

\begin{quote}
\te{జాతరలో} \te{మేము} \te{బుర్రకథ} \te{విన్నాము}.
\end{quote}

\begin{quote}
\emph{jātaralō mēmu burrakatha vinnāmu.}
\end{quote}

\begin{quote}
We heard a burrakatha at the fair.
\end{quote}

A Telugu storytelling show: a lead singer and two companions sing and talk through a tale.

\begin{cousinweb}
\textbf{\te{బుర్రకథ}} (\emph{burrakatha}) is a Telugu compound: native \emph{burra}, \textquotedblleft{}an empty gourd shell; the head\textquotedblright{} (DEDR 4332), and \emph{katha}, \textquotedblleft{}a story\textquotedblright{}, from Sanskrit.
\end{cousinweb}

\begin{grammarlens}[title={What you've built}]
One more word for telling what happened.
\end{grammarlens}

\subsection*{Your turn}
\begin{itemize}
\raggedright
  \item \emph{Say it:} \emph{burrakatha}
  \item \emph{Say it:} \emph{burrakatha}, once more
  \item \emph{Say it:} \emph{jātaralō mēmu burrakatha vinnāmu.} — We heard a burrakatha at the fair.
  \item \emph{Recall:} say the Telugu for a biography, then the Telugu for a proverb, then say \emph{burrakatha} again
\end{itemize}

\begin{tcolorbox}[breakable,colback=teal!4,colframe=teal!35!black,title={Before you move on}]
What does \te{బుర్రకథ} mean? (\textquotedblleft{}A burrakatha ballad\textquotedblright{}.) Say it once more.
\end{tcolorbox}
\section[\texorpdfstring{\te{శ్రీకారం} \te{చుట్టు} (śrīkāraṁ cuṭṭu)}{శ్రీకారం చుట్టు (śrīkāraṁ cuṭṭu)}]{\te{శ్రీకారం} \te{చుట్టు} (śrīkāraṁ cuṭṭu) — to make a start}

\label{lesson:TE-C373-srikaram-cuttu}



\begin{quote}
Before the new one: say the Telugu for a proverb, then the Telugu for a burrakatha ballad.
\end{quote}

\subsection*{You'll want to know: \te{శ్రీకారం} \te{చుట్టు}}
\textbf{\te{శ్రీకారం} \te{చుట్టు}} — \emph{śrīkāraṁ cuṭṭu} — \textquotedblleft{}to make a start\textquotedblright{}.

\begin{quote}
\te{ఆ} \te{రచయిత} \te{కొత్త} \te{నవలకు} \te{శ్రీకారం} \te{చుట్టాడు}.
\end{quote}

\begin{quote}
\emph{ā racayita kotta navalaku śrīkāraṁ cuṭṭāḍu.}
\end{quote}

\begin{quote}
The writer made a start on a new novel.
\end{quote}

Past (I …): \te{శ్రీకారం} \te{చుట్టాను}, \emph{śrīkāraṁ cuṭṭānu}.

\te{శ్రీకారం} is the auspicious letter \te{శ్రీ} written at the top of a new page; to wind it, \te{చుట్టు}, is to begin.

\begin{cousinweb}
\textbf{\te{శ్రీకారం} \te{చుట్టు}} (\emph{śrīkāraṁ cuṭṭu}) is a Telugu idiom: \emph{śrīkāraṁ}, from Sanskrit \emph{śrīkāra}, \textquotedblleft{}the word Śrī written at the top of a book or letter\textquotedblright{} (Monier-Williams), and native \emph{cuṭṭu}, \textquotedblleft{}to wind, roll\textquotedblright{} (DEDR 2715).
\end{cousinweb}

\begin{grammarlens}[title={What you've built}]
One more word for telling what happened.
\end{grammarlens}

\subsection*{Your turn}
\begin{itemize}
\raggedright
  \item \emph{Say it:} \emph{śrīkāraṁ cuṭṭu}
  \item \emph{Say it:} \emph{śrīkāraṁ cuṭṭu}, once more
  \item \emph{Say it:} \emph{ā racayita kotta navalaku śrīkāraṁ cuṭṭāḍu.} — The writer made a start on a new novel.
  \item \emph{Recall:} say the Telugu for a proverb, then the Telugu for a burrakatha ballad, then say \emph{śrīkāraṁ cuṭṭu} again
\end{itemize}

\begin{tcolorbox}[breakable,colback=teal!4,colframe=teal!35!black,title={Before you move on}]
What does \te{శ్రీకారం} \te{చుట్టు} mean? (\textquotedblleft{}To make a start\textquotedblright{}.) Say it once more.
\end{tcolorbox}
\section[(dialogue)]{Words from chapters 366–373 — a review}

\label{lesson:TE-R373-review-366-373}



\begin{quote}
Say the Telugu for a treasury and a treasure.
\end{quote}

\subsection*{Your turn}
Say these aloud:

\begin{itemize}
\raggedright
  \item a treasury, a treasure, a quest, a lord, to occupy
  \item ruins, ancient, excavations, remains, to unearth
  \item an inscription, a palm-leaf manuscript, a historian, a civilisation, an era
  \item a palace, a chariot, heritage, to demolish, to build
  \item a legend, a tale, fame, to extol, to chastise
  \item a sorcerer, magic, a ghost, a boon, a curse
  \item a storyteller, a plot, a character, imagination, a made-up story
  \item an autobiography, a biography, a proverb, a burrakatha ballad, to make a start
\end{itemize}

\begin{tcolorbox}[breakable,colback=teal!4,colframe=teal!35!black,title={Before you move on}]
A treasury? (\textbf{\te{ఖజానా}}, \emph{khajānā}.) A treasure? (\textbf{\te{నిధి}}, \emph{nidhi}.) A quest? (\textbf{\te{అన్వేషణ}}, \emph{anvēṣaṇa}.) A lord? (\textbf{\te{దొర}}, \emph{dora}.) To occupy? (\textbf{\te{ఆక్రమించు}}, \emph{ākramiñcu}.) Ruins? (\textbf{\te{శిథిలాలు}}, \emph{śithilālu}.) Ancient? (\textbf{\te{పురాతన}}, \emph{purātana}.) Excavations? (\textbf{\te{తవ్వకాలు}}, \emph{tavvakālu}.) Remains? (\textbf{\te{అవశేషాలు}}, \emph{avaśēṣālu}.) To unearth? (\textbf{\te{వెలికితీయు}}, \emph{velikitīyu}.) An inscription? (\textbf{\te{శాసనం}}, \emph{śāsanaṁ}.) A palm-leaf manuscript? (\textbf{\te{తాళపత్రం}}, \emph{tāḷapatraṁ}.) A historian? (\textbf{\te{చరిత్రకారుడు}}, \emph{caritrakāruḍu}.) A civilisation? (\textbf{\te{నాగరికత}}, \emph{nāgarikata}.) An era? (\textbf{\te{యుగం}}, \emph{yugaṁ}.) A palace? (\textbf{\te{రాజభవనం}}, \emph{rājabhavanaṁ}.) A chariot? (\textbf{\te{రథం}}, \emph{rathaṁ}.) Heritage? (\textbf{\te{వారసత్వం}}, \emph{vārasatvaṁ}.) To demolish? (\textbf{\te{కూల్చు}}, \emph{kūlcu}.) To build? (\textbf{\te{నిర్మించు}}, \emph{nirmiñcu}.) A legend? (\textbf{\te{ఐతిహ్యం}}, \emph{aitihyaṁ}.) A tale? (\textbf{\te{గాథ}}, \emph{gātha}.) Fame? (\textbf{\te{కీర్తి}}, \emph{kīrti}.) To extol? (\textbf{\te{కొనియాడు}}, \emph{koniyāḍu}.) To chastise? (\textbf{\te{దండించు}}, \emph{daṇḍiñcu}.) A sorcerer? (\textbf{\te{మాంత్రికుడు}}, \emph{māntrikuḍu}.) Magic? (\textbf{\te{మాయ}}, \emph{māya}.) A ghost? (\textbf{\te{దెయ్యం}}, \emph{deyyaṁ}.) A boon? (\textbf{\te{వరం}}, \emph{varaṁ}.) A curse? (\textbf{\te{శాపం}}, \emph{śāpaṁ}.) A storyteller? (\textbf{\te{కథకుడు}}, \emph{kathakuḍu}.) A plot? (\textbf{\te{కథాంశం}}, \emph{kathāṁśaṁ}.) A character? (\textbf{\te{పాత్ర}}, \emph{pātra}.) Imagination? (\textbf{\te{కల్పన}}, \emph{kalpana}.) A made-up story? (\textbf{\te{కట్టుకథ}}, \emph{kaṭṭukatha}.) An autobiography? (\textbf{\te{ఆత్మకథ}}, \emph{ātmakatha}.) A biography? (\textbf{\te{జీవితచరిత్ర}}, \emph{jīvitacaritra}.) A proverb? (\textbf{\te{సామెత}}, \emph{sāmeta}.) A burrakatha ballad? (\textbf{\te{బుర్రకథ}}, \emph{burrakatha}.) To make a start? (\textbf{\te{శ్రీకారం} \te{చుట్టు}}, \emph{śrīkāraṁ cuṭṭu}.)
\end{tcolorbox}
